Gravitational lensing can result in multiple images of a background source
if the source, the lensing object and the observer are nearly aligned. These
multiple images take different times to reach the observer, and if the
background source is variable, each image shows the same feature in its
variability after specific time delays. These time delay measurements are
extremely useful to estimate the current expansion rate of the Universe, the
Hubble constant.

We will consider the gravitational lens system shown in the above figure.
The left hand panel shows a galaxy cluster (lens) together with 4 images of
a background quasar formed due to gravitational lensing. The 4 images,
labelled A, B, C and D, have different fluxes as each image is magnified by
a different amount. For any given image the magnification does not change
with time. Light takes the largest time to travel for the image labelled D.

% FIGURE NEEDED: two-panel figure -- left: galaxy cluster (lens) with the
% four quasar images labelled A-D; right: light curve of image D

The light coming from this quasar is variable, and astronomers have
monitored this system for more than a decade now. The right hand panel of
the figure shows the light curve for the image D.

On the screen opposite to your station you will see a movie of the
gravitational lens system. This movie is 28 s long and will loop 6 times
with breaks of 1 minute or 2 minutes between runs. Every second on the watch
corresponds to 250.0 days in the actual lens system.

\begin{description}
\item[(OT02.1)] Let the time delays of image D with respect to images A, B
  and C be given as $t_{DA} = t_D - t_A$, $t_{DB} = t_D - t_B$, and
  $t_{DC} = t_D - t_C$, respectively. Find these time delays, taking any
  necessary steps to reduce the uncertainty in your results. \hfill[25]
\end{description}
